\section{Introduction}

Linear probes on model activations are a leading proposal for monitoring language models for deception~\citep{goldowskydill2025detecting,kretschmar2025liarsbench}. The property such a probe is supposed to detect is a mismatch between what the model believes and what it says. That is different from the statement being false. A model can say something false because it does not know the answer (a hallucination), or because it knows the answer and misreports it under some incentive (a lie). The literature draws this distinction clearly at the behavioural level~\citep{ren2025mask,pacchiardi2023catch,evans2021truthful}. Detector evaluations mostly do not: the positive class is a set of lies and the negative class is a set of honest, usually correct, answers, typically produced under a different prompt.

Such an evaluation cannot say what the detector detects. A probe that fires on any false statement, on low-knowledge states, or on the pressure prompt itself would score as well as one that fires on belief--statement mismatch. The difference matters in deployment: the first three flag honest mistakes and honest answers in adversarial-sounding contexts, and may miss a confident, knowing lie.

We study this with a matched design on a single model. Every question is asked under a neutral prompt, under pressure prompts that give a reason to misreport, and under matched control prompts with the same scenario and a truth-aligned incentive. Neutral elicitation with sampling labels each question as known or unknown for the model. Each response then falls into a cell: honest correct, hallucination, lie, honest answer under pressure, falsehood under pressure without knowledge, and feigned or genuine ignorance. We ask two questions.

\emph{How often is a false output a lie?} For Llama-3.1-8B and Gemma-2-9B on 6{,}000 TriviaQA questions, the answer depends almost entirely on how explicit the cue is. With an instruction to lie, 46--60\% of false outputs are lies. With incentive-only prompts (a bluffing game, a rival team, a deployment threat) lies are 3--15\% of false outputs, close to the 1.5--5\% that arise under control prompts, and on known questions the rate of wrong answers is at most 5.9\%. What the deployment-threat prompt does elicit, on 15--24\% of known questions, is a false claim of ignorance.

\emph{What do detectors detect?} We score all responses with an instructed-pairs deception probe~\citep{goldowskydill2025detecting}, a truth probe~\citep{marks2023geometry}, a hallucination probe, an in-domain lie probe and token-level uncertainty. Our findings are:
\begin{itemize}\itemsep1pt
\item \textbf{The usual evaluation measures the prompt.} Deception probes reach AUROC 1.00 on lies versus honest control answers and the same 1.00 on honest answers under the pressure prompt versus honest answers under the control prompt. With the prompt held fixed, the off-the-shelf deception probe separates lies from honest answers at 0.64--0.78.
\item \textbf{Within a prompt, detectors respond to falsehood in general.} A truth probe (0.84--0.88) and a hallucination probe trained only on honest errors (0.73--0.84) separate instructed and role-play lies from honest answers as well as or better than the deception probes, and the deception probes rank honest hallucinations above correct answers.
\item \textbf{Lies and knowledge-free falsehoods are confused.} No deception or truth probe separates a lie from a false answer given under the same prompt without knowledge at better than 0.79, and none separates feigned from genuine ignorance at better than 0.62. Hallucination probes treat a lie like a correct answer relative to an honest error.
\item \textbf{Incentive-only lies look like slips.} Wrong answers produced under the bluffing game are scored by uncertainty-type detectors exactly as ordinary sampling slips are.
\item \textbf{The distinction is decodable.} Probes trained with same-prompt honest answers and same-prompt honest errors as negatives reach 0.85--0.93 on lies versus all non-lies, transfer from instructed to role-play and game-incentive lies, and for instructed and role-play lies do not respond to hallucinations.
\end{itemize}

The practical implication is that a lie detector should be evaluated, and probably trained, against three kinds of negatives: honest answers under the same prompt, false answers given without knowledge under the same prompt, and honest errors under neutral prompts. An AUROC against honest correct answers under a different prompt says little about any of them. Code, prompts, graded responses and all results are in the accompanying repository.
